\section{Заключение}

В выпускной квалификационной работе рассмотрена предметная область совместного планирования поездок и обоснована необходимость интеграции ключевых пользовательских сценариев в едином мобильном приложении.

В первой главе выполнен анализ существующих решений и выявлен функциональный разрыв между специализированными сервисами и потребностью пользователей в сквозном процессе планирования. Во второй главе сформированы проектные решения для системы: пользовательские сценарии, архитектурная модель, разделение на подсистемы, модель данных и нефункциональные требования. В третьей главе описана программная реализация модулей системы и представлены результаты проверки сквозных сценариев.

Полученные результаты подтверждают, что текущая версия приложения обеспечивает целевой пользовательский процесс: управление поездкой и участниками, работу с идеями и маршрутом, учет расходов и баланса, использование погодного и AI-контекста, а также поддержку офлайн-режима с последующей синхронизацией.

Дальнейшее развитие работы связано с расширением автоматизированного тестирования, развитием аналитики пользовательских сценариев и усилением эксплуатационной устойчивости при росте нагрузки.
